\begin{tikzpicture}[
  font=\scriptsize,
  panah/.style={-{Stealth[length=2mm]}, thick},
  ukur/.style={{Stealth[length=1.6mm]}-{Stealth[length=1.6mm]}, thick},
  titik/.style={circle, draw, thick, minimum size=0.32cm, inner sep=0pt},
  kotak/.style={draw, rounded corners, align=center, text width=4.4cm,
    font=\tiny, inner sep=3pt}
]
  % Garis waktu.
  \draw[panah] (-0.4,0) -- (10.9,0) node[right, font=\tiny] {waktu};

  % Tiga titik penting.
  \node[titik, fill=blue!8] (pulih) at (2.0,0) {};
  \node[titik, fill=red!22] (bencana) at (5.0,0) {};
  \node[titik, fill=green!10] (normal) at (9.8,0) {};
  \node[align=center, anchor=north, font=\tiny] at (2.0,-0.3)
    {titik pemulihan terakhir\\yang bersih\\\textbf{contoh: 14.15}};
  \node[align=center, anchor=north, font=\tiny] at (5.0,-0.3)
    {bencana terjadi\\\textbf{contoh: 14.20}};
  \node[align=center, anchor=north, font=\tiny] at (9.8,-0.3)
    {layanan kembali\\melayani\\\textbf{contoh: 15.05}};

  % Ukuran RPO dan RTO.
  \draw[ukur, praditagreen] (2.0,0.75) -- (5.0,0.75);
  \draw[thin, gray] (2.0,0.2) -- (2.0,0.95);
  \draw[thin, gray] (5.0,0.2) -- (5.0,0.95);
  \draw[thin, gray] (9.8,0.2) -- (9.8,0.95);
  \node[font=\scriptsize\bfseries, praditagreen, anchor=south] at (3.5,0.8)
    {RPO: data yang hilang};
  \draw[ukur, red!60!black] (5.0,0.75) -- (9.8,0.75);
  \node[font=\scriptsize\bfseries, red!60!black, anchor=south] at (7.4,0.8)
    {RTO: waktu layanan berhenti};
  \node[font=\tiny, anchor=north] at (3.5,0.7) {contoh: 5 menit};
  \node[font=\tiny, anchor=north] at (7.4,0.7) {contoh: 45 menit};

  % Pertanyaan bisnis di balik kedua angka.
  \node[kotak, fill=blue!8] at (2.9,-2.05)
    {\textbf{Pertanyaan RPO:} berapa menit data terakhir boleh hilang?\\
    Ditentukan oleh seberapa sering data dicadangkan atau direplikasi.};
  \node[kotak, fill=orange!12] at (8.1,-2.05)
    {\textbf{Pertanyaan RTO:} berapa lama layanan boleh berhenti?\\
    Ditentukan oleh strategi pemulihan dan kesiapan lokasi kedua.};
\end{tikzpicture}
